\documentclass[a4paper,11pt,fleqn]{article}%fleqn pour aligner les equations à gauche
\usepackage{graphicx}
\usepackage[utf8]{inputenc}  
\usepackage[T1]{fontenc}
\usepackage{lscape}

\usepackage{tgheros,textcomp}%Police Arial
\renewcommand{\familydefault}{\sfdefault}

%\setlength{\mathindent}{0pt}%indentation dans math = 0

\usepackage[top = 1cm, bottom = 2cm, left = 1cm, right = 1cm]{geometry}
\usepackage{titlesec}
%\setcounter{secnumdepth}{4}
\usepackage{xcolor}
\usepackage{amssymb} %double flèche de réaction
\usepackage{mathrsfs}%farraday
\usepackage{xtab}%tableau sur plusieurs pages
\usepackage{esvect}%grands vecteurs

\usepackage[version=4]{mhchem} %equations chimiques

\usepackage{array}
\usepackage{chemfig}
\usepackage{multirow}
%\usepackage{sectsty}

\usepackage{caption}
\usepackage{adjustbox}
\usepackage[labelformat=simple]{subfig}
\usepackage{float}


\usepackage{amsmath}
\usepackage{amssymb}
\usepackage[cal=boondoxo]{mathalfa} 
\usepackage{enumitem}
\usepackage{lastpage}
\usepackage{tcolorbox}
\usepackage{siunitx}
\usepackage[european, straightvoltages, RPvoltages]{circuitikz}
\usepackage{tikz}
\usetikzlibrary{shapes.geometric}
\usetikzlibrary{calc}
\usetikzlibrary{automata,positioning}
\usepackage{multicol}
\usepackage{eqnarray}
\usepackage{tabularx,colortbl}%colorer fond cellules

\sisetup{
	detect-all,
	output-decimal-marker={,},
	group-minimum-digits = 3,
	group-separator={~},
	number-unit-separator={~},
	inter-unit-product={~}
} %setup pour SIunitx

\usepackage{listings}%pour insérer le code python
\definecolor{codegreen}{rgb}{0,0.6,0}
\definecolor{codegray}{rgb}{0.5,0.5,0.5}
\definecolor{codepurple}{rgb}{0.58,0,0.82}
\definecolor{backcolour}{rgb}{0.95,0.95,0.92}


\titleformat{\section}
{\titlerule
\vspace{.8ex}%
\normalfont \bfseries}
{\thesection :}{.5em}{\bfseries}

\renewcommand{\thesection}{Exercice \arabic{section}}
\titlespacing{\section}{0pt}{1em}{1em}
\renewcommand{\thesubsection}{\hspace{1cm}\alph{subsection}.}
\renewcommand\thesubfigure{\textbf{Document \arabic{subfigure} :}}

\title{\vspace{-1cm}\large 
	\begin{tabular}{|>{\centering}m{5cm}|>{\centering}m{8cm}|>{\centering\arraybackslash}m{4cm}|}
		\hline
		Correction d'exercices & \textbf{\textsc{Méthode par titrage}} &  Chapitre 3 \\
		\hline
\end{tabular}}
\date{}

\newpagestyle{mystyle}{%
	\setfoot{ROSALIE}{\thepage\,$|$\,\pageref{LastPage}}{}
}
\renewpagestyle{plain}{%
	\setfoot{ROSALIE}{\thepage\,$|$\,\pageref{LastPage}}{}
}

\pagestyle{mystyle}

\newenvironment{itemize*}%
{\begin{itemize}[label=$-$]%
		\setlength{\itemsep}{0pt}%
		\setlength{\parskip}{0pt}}%
	{\end{itemize}}

\newenvironment{itemize**}%
{\begin{itemize}[label=,labelindent=0pt, leftmargin=*,topsep=0pt]%
		\setlength{\itemsep}{0pt}%
		\setlength{\parskip}{0pt}}%
	{\end{itemize}}

\renewcommand{\arraystretch}{1.5}


\begin{document}
\maketitle
\vspace{-2.5cm}
\begin{multicols}{2}

\section*{Exercice 4 : Dans une eau de brassage}
Le titrage donne un volume équivalent $V_E$ = \SI{12}{mL}.\\
L'équation support du dosage permet d'écrire qu'à l'équivalence $n(\ce{Ag+}) = n_i(\ce{Cl-}) $. Ainsi on a : $$c_2 \cdot V_E = c \cdot V_0$$
soit  $$c = \dfrac{c_2 \cdot V_E}{ V_0}$$
\fbox{AN : } $c= \dfrac{ \num{10E-3} \times \num{12E-3}}{ \num{100E-3}} = \SI{1.2E-3}{\mole\per\liter}$

La concentration en ions \ce{Cl-} de l'échantillon est de \SI{1.2E-3}{\mole\per\liter}. Elle correspond à une concentration en masse de : $$Cm = C \cdot M$$
\fbox{AN : } $Cm=\num{1.2E-3} \times \num{35.5} = \SI{4.2E-2}{\gram\per\liter}$

Soit une concentration en ions \ce{Cl-} de \SI{42}{\milli\gram\per\liter}. En prenant en compte les données de l'énoncé, elle ne conviendrait pas à une bière brune car trop peu concentrée en ions \ce{Cl-}.\\

\section*{Exercice 5 : Titrage de l'ibuprofène}
\begin{enumerate}
	\item La masse maximale d'ibuprofène soluble dans \SI{40}{mL} d'eau s'obtient en multipliant la solubilité de la molécule par ce volume : $$m_{max} = s \cdot V$$
	\fbox{AN : } $m_{max} = \num{0.043} \times 40 = \SI{1.72}{mg}$\\
	La masse est très faible c'est la raison pour laquelle on préfèrera utiliser l'éthanol pour solubiliser d'avantage d'ibuprofène.
	\item Equation support : $$\ce{R-COOH(aq) + HO_(aq) -> R-COO-(aq)  + H2O(l)}$$
	\item Courbe 1 : $pH=f(V_B)$ \\ Courbe 2 : $\dfrac{dpH}{dV_B = f(V_B)}$
	\item L'équivalence correspond à l'instant du titrage où les quantités de matières de réactifs titrant et titrés sont introduits dans des proportions st\oe chimoétriques. On a alors $$n(\ce{RCOOH}) = n_{eq}(\ce{HO-})$$
	\item On peut lire graphiquement un volume équivalent $V_E$ = \SI{9.75}{mL}
	\item Quantité de matière d'ibuprofène dans la solution titrée :
	$$n(\ce{RCOOH}) = n_{eq}(\ce{HO-})$$
	$$n(\ce{RCOOH}) =c_B \cdot V_E$$
\fbox{AN : } $n(\ce{RCOOH}) = \num{0.20} \times \num{9.75E-3} = \SI{1.95E-3}{\mol}$
Masse d'ibuprofène correspondante :
$$m=n \cdot M$$
$$m(\ce{RCOOH}) = n(\ce{RCOOH}) \cdot M(\ce{RCOOH})$$
\fbox{AN : } $m(\ce{RCOOH}) =\num{1.95E-3} \times \num{206.0} =\SI{0.402}{g} $

La masse d'ibuprofène contenu dans le comprimé est de \SI{402}{mg}.
	
\end{enumerate}
\end{multicols}
\end{document}